\documentclass[12pt]{article}
\usepackage[utf8]{inputenc}
\usepackage[T1]{fontenc}
\usepackage{hyperref}
\title{Zamwal Language Resources for K-12 Curriculum Development}
\author{Compiled by Claude Sonnet 5.5 (claude-sonnet-5-5) for Andr\'as Kornai}
\date{}
\begin{document}
\maketitle

\section{Introduction}
Zamwal (ISO 639-3 zah, Glottocode zang1255) is an endangered Chadic language of the Afro-Asiatic family, spoken in Bauchi State, Nigeria. In the sources I found it appears mostly as Zangwal, also Zwangal or Twar. About 100 speakers were reported in 1993. The reference description says only older adults use it as a first language and that it is not known to be taught in schools. The searches turned up no teaching materials of any kind.

\section{Primary Educational Resources}
I found no primers, textbooks, readers or school programmes for this language.

\section{Linguistic and Reference Materials}
\begin{itemize}
\item Wikipedia, ``Zangwal language'': an overview of the language and its classification. \url{https://en.wikipedia.org/wiki/Zangwal_language}
\item Glottolog entry (zang1255): classification and references. \url{https://glottolog.org/resource/languoid/id/zang1255.bigmap.html}
\item Endangered Languages Project entry, with a bibliography. \url{https://endangeredlanguages.com/lang/3146?hl=en}
\item Wikipedia, ``Chadic languages'': background on the wider family. \url{https://en.wikipedia.org/wiki/Chadic_languages}
\end{itemize}
Search results also pointed to a paper by Blench (2006) on the Shall and Zwall languages of Central Nigeria. I did not see its contents and give no link for it.

\section{Religious and Community Resources}
None found.

\section{Digital and Audio Resources}
No audio, video or software was found.

\section{Curriculum Development Context}
Very little is documented and the speaker base is very small and elderly. Any curriculum would first need documentation work: recording the remaining speakers and describing the orthography, with the community's involvement. Nothing in the search results shows that a standard orthography exists. Nigerian national languages, English and Hausa, would carry instruction in the meantime.
\end{document}